\section{Open problems and proposed experiments}\label{sec:open}

% Survey section (draft 2026-10-15). Each problem is phrased as the
% measurement or proof that would settle it, per the repository's
% "name the obstruction" rule.

\paragraph{1. The solving degree of Weil-descended summation systems.}
Every subexponential claim for prime-degree binary curves rests on a bound
for this degree (\S\ref{sec:theory}); none is proved and the experimental
evidence is contested. \emph{Experiment:} measure first-fall, last-fall
and solving degree of $\Sem{m+1}$ restricted to $\ell$-dimensional
subspaces over a grid of $(n,\ell,m)$ with $m\ge3$, at least four values
of $n$ per $(\ell,m)$, against the semi-regular prediction, with
multi-graded regularity (one block per summand) as a second model.
\emph{Settles:} whether the degree grows with $n$ at fixed $(\ell,m)$; a
growth consistent with zero would revive the Petit--Quisquater and Semaev
heuristics, a positive slope would close them for this model.

\paragraph{2. Cold index calculus against rho with every operation
priced.} No measured arm beats matched rho on the cold one-target boundary
once native work is priced on both sides (\S\ref{sec:frontier}).
\emph{Experiment:} the frontier campaign of \S\ref{sec:frontier} extended
to the largest sizes each arm can reach, with exponents fitted per arm and
the crossover reported only where the fitted exponents separate across the
whole pricing band. \emph{Settles:} whether any implemented arm has an
exponent below $1/2$ in $r$ on these families.

\paragraph{3. Relation yield is now the dominant stage.} Root finding made
the decomposition oracle cheap; what remains is the number of trials per
relation, governed by the factor-base density. \emph{Experiment:} for each
factor-base family, the measured trials-per-relation law against $r$ and
base size, compared with the counting prediction $m!\,r/N^{m-1}$; the
residual (e.g.\ the unexplained factor 5 in the ECC2K-130 extrapolation,
\S\ref{sec:koblitz}) is the target.

\paragraph{4. Root finding beyond two summands on prime curves.} The
two-summand oracle is prior art in mechanism (McGuire--Mueller) and costs
$\Theta(N)$ per trial; three summands cost a pair table of $N^2$ entries.
\emph{Question:} is there a sub-$N^2$ structure for $\Sem4$ membership
over a prime-field base---for example from the $D_m$-invariant coordinates
of 2-torsion curves---and does it change the $r^{2/3}$ cost of two-summand
prime index calculus? A negative answer with a counting argument would be
as useful as a positive one.

\paragraph{5. The preprocessing regime.} Index calculus with a retained
table beats rho without one, and loses by $53$--$317\times$ to
Bernstein--Lange rho with an equal or smaller table (\S\ref{sec:ladder}).
\emph{Experiment:} measure the advice--time curve $(S,T)$ of the
compact-orbit arm across sizes against preprocessing rho's
$S\,T^2\approx N$; an index-calculus online result counts only if it lies
below that curve at equal $S$.

\paragraph{6. Symmetry at $k=4$, end to end.} The $D_4$ reduction and its
four-unknown weighted form move a \emph{modelled} handover over $\FF_{p^4}$
from about $2^{150}$ to $2^{119}$--$2^{137}$ depending on frame
(\S\ref{sec:ladder}).
\emph{Experiment:} a complete $k=4$ index calculus over $\FF_{p^4}$ at
several $p$, cold, priced, against rho with the available automorphisms,
to replace the model by a measured slope.

\paragraph{7. Pricing and calibration infrastructure.} Rankings now
require every arm's native work to be counted (AGENTS.md \S5a).
\emph{Work:} count the three uncounted categories of the shared-rank arm;
calibrate every ladder size on one fixed host; resolve the inherited-F4
cross-host counter drift; register the $n=83$ dense root index in the
harness.

\paragraph{8. Genus three after cubic GHS, to the crossover.} The one
family whose fitted exponent ($q^{1.380}$) lies below rho's $q^{3/2}$ is
still $1.7$--$2.6\times$ the negation-map rho model at $q=2^{20}$, and its
local exponent rises by $0.02$ per bit because the linear algebra grows as
$q^{1.5}$ (\S\ref{sec:ladder}).
\emph{Experiment:} rebalance the factor-base fraction to the structured
linear algebra, price a multiplication mod $\ell$ at its measured cost, and
run to $q\ge2^{22}$ against a measured negation-map rho on the same
instances, in the harness and across the pricing band. A measured
crossover there would be the programme's first \cls{advance} against rho,
on a narrow and undeployed family.
